\chapter{Practical Usage Examples}\label{ch:usage}

The following scenarios were executed on the full dataset (23.76 million rows) during the
demonstration rehearsal. The outputs below are copied from \code{results/demo/}, with blank
lines and timing lines removed. Before step~1 the policy had been applied as of the previous
month (2026-10), so the scenarios show the routine monthly situation.

Every scenario can also be run from a web dashboard (\code{scripts/07\_dashboard.sh}, then
\url{http://localhost:8501}). It is built with Streamlit and shows the tiers, the business query
on both designs, the dry run, the incident and its recovery, the monthly job and the security
checks, each started by one button. The dashboard calls the same \code{DLM\_PKG} procedures as
the scripts, and serialises write operations so two clicks cannot run a job twice.
\code{scripts/07\_dashboard.sh --fresh} reloads the data and prepares the demonstration state.

\section{Scenario 1 -- State of the tiers}

\needspace{0.3\textheight}
\begin{lstlisting}[style=console, caption={\code{scripts/06\_demo.sh 1} (excerpt)}]
=== BILLING_RECORDS partitions by tier ===
TIER     TABLESPACE_N COMPRESS PARTITIONS FIRST_PART LAST_PART    NUM_ROWS
-------- ------------ -------- ---------- ---------- ---------- ----------
HOT      TELCO_HOT    DISABLED         13 P_202510   P_202610      5359479
WARM     TELCO_ARCH   ENABLED          24 P_202310   P_202509      8656225
=== BILLING_ARCHIVE partitions ===
PARTITIONS FIRST_PART LAST_PART    NUM_ROWS
---------- ---------- ---------- ----------
        24 P_202110   P_202309      7049916
\end{lstlisting}

Each tier is in the expected tablespace and format. Retained history is split between the
online table (hot + warm) and the archive.

\section{Scenario 2 -- Month-end revenue report (business scenario)}

Finance asks for October 2026 revenue by service. The same statement runs on both designs and
returns identical results (4 rows: 128\,397 voice, 64\,382 SMS, 214\,795 data and 21\,426
roaming lines).

\needspace{0.3\textheight}
\begin{lstlisting}[style=console, caption={Execution plans of the same query on both designs}]
-- Baseline (heap table)
| Id  | Operation          | Name                 | Rows  |
|   0 | SELECT STATEMENT   |                      |       |
|   1 |  HASH GROUP BY     |                      |     4 |
|   2 |   TABLE ACCESS FULL| BILLING_RECORDS_BASE |   336K|

-- Partitioned design
| Id  | Operation               | Name            | Rows  | Pstart| Pstop |
|   0 | SELECT STATEMENT        |                 |       |       |       |
|   1 |  HASH GROUP BY          |                 |     4 |       |       |
|   2 |   PARTITION RANGE SINGLE|                 |   429K|    73 |    73 |
|   3 |    TABLE ACCESS FULL    | BILLING_RECORDS |   429K|    73 |    73 |
\end{lstlisting}

The baseline must read the whole six-year table. One month is about 1.8\,\% of the rows,
still too many for the \code{bill\_date} index. The partitioned plan reads one partition. For
interval partitioning, Pstart/Pstop is the interval number counted from the first partition
(month 73), not the partition's current position.

\section{Scenario 3 -- Dry run of the monthly job}

\needspace{0.3\textheight}
\begin{lstlisting}[style=console, caption={\code{exec dlm\_pkg.run\_lifecycle(p\_dry\_run => true)} for 2026-11}]
ACTION       MONTHS FIRST_MONTH   LAST_MONTH
-------- ---------- ------------- ------------
PURGE             1 2021-10       2021-10
WARM              1 2025-10       2025-10
ARCHIVE           1 2023-10       2023-10
\end{lstlisting}

The operator can check exactly what the job will do before anything changes.

\section{Scenario 4 -- Incident: archive job crashes, then recovers}

We simulate a crash right after the first exchange of the archive of 2023-10
(\code{p\_simulate\_failure => 'AFTER\_STEP1'}).

\needspace{0.3\textheight}
\begin{lstlisting}[style=console, caption={Incident and recovery (\code{scripts/06\_demo.sh 4}, 16 s in total)}]
=== Rows of 2023-10 before the incident ===
     328606
=== Archive job crashes right after step 1 (exchange out) ===
ORA-20020: Simulated failure after step 1
=== Where is the data? Not lost: it waits in BILLING_XCHG ===
ROWS_IN_EXCHANGE_TABLE      328606
=== Integrity check (2021-10 .. 2023-10) detects the incomplete month ===
MONTHS_IN_ERROR             1
=== Recovery: simply re-run the job ===
PL/SQL procedure successfully completed.
ROWS_ARCHIVED               328606
ROWS_IN_EXCHANGE_TABLE      0
MONTHS_IN_ERROR             0
=== Audit trail ===
ACTION    TARGET       STATUS   MESSAGE
VALIDATE  BILLING_ALL  OK       25/25 months match control totals (from 2021-10 to 2023-10)
ARCHIVE   P_202310     OK       resumed after interruption; exchanged into BILLING_ARCHIVE, read only
ARCHIVE   2023-10      RESUMED  found staged rows from an interrupted run - resuming at step 2
VALIDATE  BILLING_ALL  ERROR    24/25 months match control totals (from 2021-10 to 2023-10)
ARCHIVE   P_202310     ERROR    SIMULATED FAILURE after step 1 (data is safe in BILLING_XCHG)
\end{lstlisting}

No row was lost. The integrity check detected the incomplete month, a simple re-run resumed
at step~2, and the audit trail records each event.

\section{Scenario 5 -- Guard and monthly job}

\needspace{0.3\textheight}
\begin{lstlisting}[style=console, caption={Guard, then the monthly job for 2026-11 (46 s)}]
SQL> exec dlm_pkg.purge_partition(date '2022-01-01', p_retention_months => 24)
ORA-20010: PURGE 2022-01 refused: requested retention 24 months is below the
           legal floor of 60 months

SQL> exec dlm_pkg.run_lifecycle
ACTION     STATUS     OPERATIONS  TOTAL_SEC ROWS_PROCESSED
---------- ---------- ---------- ---------- --------------
PURGE      OK                  1        .07         261676
WARM       OK                  1       3.22         395535
VALIDATE   OK                  1      15.57
RUN_END    OK                  1      45.51
MESSAGE: 48/48 months match control totals (from 2021-11 to 2025-10)

=== Archived history is still one query away (BILLING_ALL) ===
DISPUTED_2022_03        2740
\end{lstlisting}

Purging a month of 261\,676 rows took 0.07~s. Compressing one month took 3.2~s. Most of the
46-second run is statistics and the validation of the 48 frozen months (16~s).

\section{Scenario 6 -- What the billing application cannot do}

\needspace{0.3\textheight}
\begin{lstlisting}[style=console, caption={Asserted security checks as \code{TELCO\_APP} (excerpt)}]
PASS  S02 insert a bill in a HOT month (rolled back)
PASS  S03 late insert dated in an archived month (2022-03) rejected  (ORA-1402)
PASS  S05 direct write into the base table rejected  (ORA-942)
PASS  S06 delete billing records rejected  (ORA-1031)
PASS  S09 change retention policy rejected  (ORA-942)
PASS  S10 run a lifecycle job rejected  (ORA-6550)
PASS  S14 lock of reference data rejected (READ grant)  (ORA-1031)
Security checks (TELCO_APP) passed: 14   failed: 0
\end{lstlisting}
